% Deep learning theory (Harshil)
% Receptive fields measured from input gradients of src/dl_models.py (random weights, one output pixel,
% all positions in a 16-pixel range). Numbers from reports/dl/results_dl.csv, reports/dl/gbm_oof_metadata.json,
% reports/dl/hybrid_ensemble_results.csv and reports/dl/architecture_ablation_summary.md.

\section{Deep Learning Theory}
\label{sec:dl_theory}

This section explains four results of Phase~2: why depth helps and width does not, why a model trained on
crops fails on full patches, why rotations do not help, and why the Phase~1 model adds nothing to the
U-Net.

\subsection{Receptive field}
\label{sec:dl_theory_rf}

An output pixel of a CNN depends only on a window of the input, its receptive field. For a stack of
layers with kernel size $k_l$ and stride $s_l$, its width is
\begin{equation}
r_L = 1 + \sum_{l=1}^{L} (k_l - 1) \prod_{m=1}^{l-1} s_m .
\label{eq:rf}
\end{equation}
Each $3\times3$ convolution adds $2\,J_l$ pixels, where $J_l = \prod_{m<l} s_m$ is the spacing of its
inputs in pixels. A convolution at U-Net level $j$, after $j-1$ halvings, therefore adds $2^{j}$ pixels, so
each extra level roughly doubles the reach. Measured from input gradients, the field is 9\,px for our
CNN (four convolutions, no downsampling), 18\,px for a U-Net with 2 levels, 44\,px with 3 levels and
92--100\,px with 4 levels; the last depends on the pixel's position relative to the pooling grid. A
patch is 64\,px wide. At depth 3 no output pixel sees the whole patch. At depth 4, every pixel within
14\,px of the centre, where the fire is, does.

The theoretical field overstates what the model uses. Luo et al.~\cite{luo2016} show that the effective
receptive field is close to Gaussian and that its width grows only as $\sqrt{n}$ for $n$ stacked
convolutions at one resolution, whereas downsampling scales it directly. Width adds parameters at every
level but no reach. This matches our ablation: at almost the same parameter count, 4 levels at width 32
(1.93\,M) beat 3 levels at width 64 (1.87\,M) by 0.065 test PR AUC in single runs, and the 3-seed
comparison of depth 4 and depth 3 gives 0.344 $\pm$ 0.007 against 0.280 $\pm$ 0.002. The gain is at
low false-alarm rates (recall at 1\% FPR 0.399 $\to$ 0.459, ROC AUC unchanged). This is where a pixel
model fails. Phase~1 missed most fire more than 5\,km from the current front. A pixel's fire
risk there depends on the shape of the front and on wind and fuel over tens of kilometres, which
fixed-window features summarise only coarsely.

\subsection{Translation equivariance and input size}

A convolution commutes with shifts: if $T_v$ moves an image by $v$ pixels, then
$f(T_v\mathbf{x}) = T_v f(\mathbf{x})$ away from the borders. The same weights are applied everywhere,
so a fully convolutional network~\cite{long2015} accepts any input size. This makes it possible to train
on $32\times32$ crops and predict on $64\times64$ patches. It does not make the result good. In our
centre-crop run (0.172 against 0.269, single runs) three things shift between training and test:
\begin{itemize}
\item the share of fire, about 3 times higher in the patch centre;
\item the zero-padded borders, which now lie in the middle of the field;
\item the statistics that batch normalisation~\cite{ioffe2015} learned on crops.
\end{itemize}
Equivariance guarantees that the computation transfers, not that the training distribution does.

\subsection{Rotations and the wind field}

Flips and $90^\circ$ rotations form a group of 8 symmetries of the square grid. A group-equivariant
CNN~\cite{cohen2016} satisfies $f(g\mathbf{x}) = g f(\mathbf{x})$ for every symmetry $g$ by
construction; augmentation only encourages it. On our data the symmetry has two problems. First, the
wind direction $\theta$ is a direction, not a pixel value: a quarter turn $k$ of the grid must also turn
the wind, $\theta \to \theta + 90^\circ k$ on our south-up grid (Section~\ref{sec:dl_aug}). Second, the
data are not symmetric: north and south differ in sun exposure and terrain aspect, and U.S.\ fires have
prevailing wind directions. A rotated patch is a plausible image but an unlikely fire. This is why flips
and rotations did not help, and why we did not build the symmetry into the architecture.

\subsection{Stacking, leakage and ensembles}

The hybrid U-Net receives the XGBoost probability $\hat p(\mathbf{x})$ as a 13th channel. On the
training patches this channel has to come from models that did not see them~\cite{wolpert1992}: with
$K$ folds and patch $i$ in fold $g(i)$,
\begin{equation}
\hat p_i = f_{-g(i)}(\mathbf{x}_i),
\label{eq:oof}
\end{equation}
where $f_{-g}$ is fitted without fold $g$. If neighbouring patches of the same fire fall into different
folds, $f_{-g(i)}$ has still seen that fire, and $\hat p_i$ looks more reliable on training data than on
test data. The U-Net then learns to over-trust it. Grouping the folds by fire event lowered the training
out-of-fold PR AUC from 0.306 to 0.268, against 0.207 on eval and 0.257 on test.

Even without leakage, the channel helped by less than 0.003. The simplest explanation is redundancy:
everything XGBoost knows comes from the same 12 channels the U-Net reads, through features computed over
windows that the depth-4 U-Net already covers. The ensemble
$p = \alpha\, p_{\text{U-Net}} + (1-\alpha)\, p_{\text{XGB}}$ says the same thing. With $\alpha$ tuned
on eval, depth 3 gets $\alpha = 0.72$ and gains 0.005 (within noise), and depth 4 gets $\alpha = 1$:
XGBoost receives no weight. Forcing $\alpha = 0.5$ lowers depth-4 PR AUC by 0.042.

\subsection{Optimisation and regularisation}

All models are trained with Adam~\cite{kingma2015} at a learning rate of $10^{-3}$ and batch size 32, with
batch normalisation after every convolution~\cite{ioffe2015}. There is no dropout or weight decay. Early
stopping~\cite{prechelt1998} is the regulariser: training stops after 5 epochs without a gain in eval PR
AUC, and the best epoch is kept (epochs 9--16 for the final model). Eval PR AUC peaks while the training
loss is still falling, so without early stopping the models overfit.
